\documentclass[aspectratio=169,10pt]{beamer}
% ---------------------------------------------------------------
% R4.02 Qualite de developpement -- preambule commun des diapositives
% Charte reprise de R3.07 (metropolis, dark tech)
% Palette IUT partagee (violet du logo + accent or)
% Utilisation :
%   \documentclass[aspectratio=169,10pt]{beamer}
%   % ---------------------------------------------------------------
% R4.02 Qualite de developpement -- preambule commun des diapositives
% Charte reprise de R3.07 (metropolis, dark tech)
% Palette IUT partagee (violet du logo + accent or)
% Utilisation :
%   \documentclass[aspectratio=169,10pt]{beamer}
%   % ---------------------------------------------------------------
% R4.02 Qualite de developpement -- preambule commun des diapositives
% Charte reprise de R3.07 (metropolis, dark tech)
% Palette IUT partagee (violet du logo + accent or)
% Utilisation :
%   \documentclass[aspectratio=169,10pt]{beamer}
%   \input{../../latex/common/r402-slides-preamble.tex}
%   \graphicspath{{../../latex/common/assets/}{img/}{figs/}}
% ---------------------------------------------------------------
\usetheme{metropolis}
\usepackage{appendixnumberbeamer}

\usepackage[T1]{fontenc}
\usepackage[utf8]{inputenc}
\usepackage[french]{babel}
\usepackage[scale=0.94]{tgheros}
\renewcommand{\familydefault}{\sfdefault}
\usepackage{amsmath,amssymb,amsfonts}
\usepackage{graphicx}
\usepackage{booktabs}
\usepackage{multirow}
\usepackage{tikz}
\usetikzlibrary{shapes.geometric, shapes, arrows, arrows.meta, positioning, calc, fit}
\usepackage{xcolor}
\usepackage{listings}
\usepackage[skins,listings]{tcolorbox}
\usepackage{qrcode}
\usepackage{eurosym}
\usepackage{hyperref}

% ---------------------------------------------------------------
% PALETTE (dark-mode tech: near-black canvas, violet (logo IUT) + gold accents)
% ---------------------------------------------------------------
\definecolor{qabg}{HTML}{0E1016}      % canvas background
\definecolor{qapanel}{HTML}{181B24}   % card / block panel
\definecolor{qapanelalt}{HTML}{20242F}% secondary panel tint
\definecolor{qaline}{HTML}{2B303D}    % hairline rules / table rules
\definecolor{qatext}{HTML}{ECEEF2}    % primary text (off-white)
\definecolor{qamuted}{HTML}{8B93A6}   % secondary / caption text
\definecolor{qaaccent}{HTML}{D93ED7}  % primary accent -- violet du logo IUT, eclairci pour fond sombre
\definecolor{qaaccent2}{HTML}{F2B84B} % secondary accent -- or/ambre
\definecolor{qacodekw}{HTML}{C792EA}  % code : mots-cles
\definecolor{qacodestr}{HTML}{C3E88D} % code : chaines
\definecolor{qacodecom}{HTML}{7F8AA3} % code : commentaires
\definecolor{qacodeann}{HTML}{F2B84B} % code : annotations / javadoc

\setbeamercolor{background canvas}{bg=qabg}
\setbeamercolor{normal text}{fg=qatext}
\setbeamercolor{palette primary}{bg=qabg, fg=qamuted}
\setbeamercolor{frametitle}{bg=qabg, fg=qatext}
\setbeamercolor{progress bar}{fg=qaaccent, bg=qaline}
\setbeamercolor{alerted text}{fg=qaaccent2}
\setbeamercolor{block title}{bg=qapanel, fg=qaaccent}
\setbeamercolor{block body}{bg=qapanel, fg=qatext}
\setbeamercolor{block title alerted}{bg=qapanel, fg=qaaccent2}
\setbeamercolor{block body alerted}{bg=qapanel, fg=qatext}
\setbeamercolor{block title example}{bg=qapanelalt, fg=qaaccent}
\setbeamercolor{block body example}{bg=qapanelalt, fg=qatext}
\setbeamercolor{example text}{fg=qaaccent}
\setbeamercolor{itemize item}{fg=qaaccent}
\setbeamercolor{itemize subitem}{fg=qaaccent2}
\setbeamercolor{enumerate item}{fg=qaaccent}
\setbeamercolor{enumerate subitem}{fg=qaaccent2}
\setbeamercolor{caption}{fg=qamuted}
\setbeamercolor{page number in head/foot}{fg=qamuted, bg=qabg}
\setbeamercolor{title}{fg=qatext}
\setbeamercolor{subtitle}{fg=qamuted}
\setbeamercolor{author}{fg=qatext}
\setbeamercolor{institute}{fg=qamuted}
\setbeamercolor{date}{fg=qamuted}
\setbeamercolor{section in toc}{fg=qaaccent}
\setbeamercolor{subsection in toc}{fg=qamuted}
% Sommaire de debut de section : section en cours en violet et gras, les autres en gris
\setbeamercolor{section in toc shaded}{fg=qamuted}
\setbeamercolor{subsection in toc shaded}{fg=qamuted}

\setbeamertemplate{blocks}[rounded][shadow=false]
\setbeamertemplate{itemize item}{\raisebox{0.15em}{\scriptsize\textcolor{qaaccent}{\ensuremath{\blacktriangleright}}}}
\setbeamertemplate{itemize subitem}{\raisebox{0.1em}{\tiny\textcolor{qaaccent2}{\ensuremath{\bullet}}}}

\setbeamerfont{frametitle}{series=\bfseries, size=\Large}
\setbeamerfont{title}{series=\bfseries, size=\Large}
\setbeamerfont{subtitle}{size=\normalsize}
\setbeamerfont{institute}{size=\scriptsize}
\setbeamerfont{date}{size=\scriptsize}
\setbeamerfont{author}{size=\small}

\metroset{sectionpage=none, subsectionpage=none, numbering=fraction, progressbar=frametitle}
\setbeamertemplate{section in toc}[sections numbered]

% ---------------------------------------------------------------
% LISTINGS -- theme sombre (lstlisting reste utilisable tel quel)
% \lstcode : police + couleur de base, a reutiliser dans les
%            surcharges locales (ex. basicstyle=\lstcode\scriptsize)
% ---------------------------------------------------------------
\newcommand{\lstcode}{\ttfamily\color{qatext}}

\lstdefinelanguage{TypeScript}{
  keywords={typeof, new, true, false, catch, function, return, null, switch, var,
    if, in, while, do, else, case, break, const, let, interface, class, extends,
    implements, enum, export, import, this, super, constructor, public, private,
    protected, static, async, await, throw, throws, try, finally, void},
  ndkeywords={boolean, any, number, string, int, long, double},
  sensitive=true,
  comment=[l]{//},
  morecomment=[s]{/*}{*/},
  morestring=[b]',
  morestring=[b]"
}

\lstset{
  language=Java,
  basicstyle=\lstcode\footnotesize,
  keywordstyle=\color{qacodekw}\bfseries,
  ndkeywordstyle=\color{qaaccent2}\bfseries,
  identifierstyle=\color{qatext},
  stringstyle=\color{qacodestr},
  commentstyle=\color{qacodecom}\itshape,
  morecomment=[s][\color{qacodeann}]{/**}{*/},
  backgroundcolor=\color{qapanelalt},
  rulecolor=\color{qaline},
  frame=single,
  framesep=3pt,
  xleftmargin=4pt,
  xrightmargin=4pt,
  keepspaces=true,
  columns=fullflexible,
  showstringspaces=false,
  breaklines=true,
  tabsize=2,
  escapeinside={(*@}{@*)},
  % accents UTF-8 dans le code (listings ne gere pas l'UTF-8 multi-octets)
  literate={é}{{\'e}}1 {è}{{\`e}}1 {ê}{{\^e}}1 {à}{{\`a}}1 {â}{{\^a}}1
           {ù}{{\`u}}1 {û}{{\^u}}1 {ô}{{\^o}}1 {î}{{\^i}}1 {ç}{{\c c}}1
           {É}{{\'E}}1 {À}{{\`A}}1 {€}{{\euro{}}}1 {…}{{\ldots}}1,
}

% ---------------------------------------------------------------
% CARD-STYLE BOXES
% ---------------------------------------------------------------
\tcbset{
  figcard/.style={colback=white, colframe=white, boxrule=0pt, arc=4pt,
    left=3pt, right=3pt, top=3pt, bottom=3pt, boxsep=0pt,
    halign=center, valign=center},
  datacard/.style={colback=qapanel, colframe=qaline, boxrule=0.4pt, arc=3pt,
    left=8pt, right=8pt, top=6pt, bottom=6pt, coltext=qatext,
    before upper={\color{qatext}}, halign=center},
  codebox/.style={colback=qapanelalt, colframe=qaline, boxrule=0.4pt, arc=3pt,
    left=8pt, right=8pt, top=5pt, bottom=5pt, listing only,
    listing options={basicstyle=\ttfamily\small\color{qaaccent},
      breaklines=true, showstringspaces=false, columns=fullflexible,
      keepspaces=true, frame=none, backgroundcolor={}}}
}
\newtcolorbox{figcard}[1][]{figcard,#1}
\newtcolorbox{datacard}{datacard}
\newtcblisting{codebox}{codebox}

% ---------------------------------------------------------------
% HANDY MACROS
% ---------------------------------------------------------------
% Mention d'elaboration (discrete : plus sombre que le texte secondaire)
\definecolor{qaaimention}{HTML}{4E5566}
\newcommand{\raimention}{{\tiny\color{qaaimention}Supports \'elabor\'es avec l'aide de Claude Code Opus 5.5}}
\newcommand{\hl}[1]{\alert{#1}}
\newcommand{\source}[1]{{\scriptsize\color{qamuted}#1}}

% Sommaire automatique en debut de section
\AtBeginSection[]{%
  \begin{frame}<beamer>{Plan du cours}
    \tableofcontents[currentsection]
  \end{frame}}

% Page de titre : \rtitleslide{numero du cours}{titre du cours}
% Les auteurs sont repris de \rauthors (modifiable avant \begin{document})
\newcommand{\rauthors}{Chouki Tibermacine}
\newcommand{\rtitleslide}[2]{%
\begin{frame}[plain,noframenumbering]
  \begin{center}
    \vspace{2mm}
    {\color{qaaccent}\rule{2.1cm}{1.2pt}}\\[6mm]
    {\Large\bfseries\color{qatext} R4.02}\\[1mm]
    {\Large\bfseries\color{qatext} Qualit\'e de d\'eveloppement}\\[5mm]
    {\large\color{qamuted} Cours #1}\\[1mm]
    {\large\color{qaaccent} #2}\\[6mm]
    \includegraphics[height=2.1cm]{iut_logo_transparent.png}

    \vspace{4mm}

    {\small\color{qatext} \rauthors}\\[1mm]
    {\scriptsize\color{qamuted} BUT 2 Informatique --- IUT de Vannes --- Universit\'e Bretagne Sud}\\[2mm]
    \raimention
  \end{center}
\end{frame}
}

% ---- Diapo de demonstration : \rdemo{numero}{titre}{contenu (items)}{fichier(s) a ouvrir} ----
\newcommand{\rdemo}[4]{%
\begin{frame}{D\'emo #1 --- #2}
  \begin{tikzpicture}[remember picture, overlay]
    \node[anchor=north east, fill=qaaccent2, text=qabg, font=\scriptsize\bfseries,
      rounded corners=3pt, inner sep=4pt] at ([xshift=-0.4cm, yshift=-0.35cm]current page.north east) {D\'EMO};
  \end{tikzpicture}
  \begin{columns}[T]
    \begin{column}{0.62\textwidth}
      \begin{itemize}
        #3
      \end{itemize}
    \end{column}
    \begin{column}{0.34\textwidth}
      \begin{datacard}
        \raggedright\scriptsize
        {\color{qamuted} Dans l'IDE}\\[1mm]
        #4
      \end{datacard}
    \end{column}
  \end{columns}
\end{frame}}

% ---- Diapo de fin « Questions ? » ----
\newcommand{\rcontact}{chouki.tibermacine@univ-ubs.fr}
\newcommand{\rquestionslide}{%
\begin{frame}{Questions ?}
  \begin{center}
    \vspace{2mm}
    \begin{tikzpicture}
      \node[regular polygon, regular polygon sides=6, draw=qaline, thick,
        minimum size=4.6cm, rotate=90] (hexin) {};
      \node[regular polygon, regular polygon sides=6, draw=qaaccent, dashed,
        thick, minimum size=6.2cm, rotate=90] (hexout) {};
      \node[font=\fontsize{92}{92}\selectfont\bfseries, color=qaaccent2] at (hexin) {?};
    \end{tikzpicture}\\[5mm]
    {\small\color{qamuted} \rcontact}
  \end{center}
\end{frame}}

%   \graphicspath{{../../latex/common/assets/}{img/}{figs/}}
% ---------------------------------------------------------------
\usetheme{metropolis}
\usepackage{appendixnumberbeamer}

\usepackage[T1]{fontenc}
\usepackage[utf8]{inputenc}
\usepackage[french]{babel}
\usepackage[scale=0.94]{tgheros}
\renewcommand{\familydefault}{\sfdefault}
\usepackage{amsmath,amssymb,amsfonts}
\usepackage{graphicx}
\usepackage{booktabs}
\usepackage{multirow}
\usepackage{tikz}
\usetikzlibrary{shapes.geometric, shapes, arrows, arrows.meta, positioning, calc, fit}
\usepackage{xcolor}
\usepackage{listings}
\usepackage[skins,listings]{tcolorbox}
\usepackage{qrcode}
\usepackage{eurosym}
\usepackage{hyperref}

% ---------------------------------------------------------------
% PALETTE (dark-mode tech: near-black canvas, violet (logo IUT) + gold accents)
% ---------------------------------------------------------------
\definecolor{qabg}{HTML}{0E1016}      % canvas background
\definecolor{qapanel}{HTML}{181B24}   % card / block panel
\definecolor{qapanelalt}{HTML}{20242F}% secondary panel tint
\definecolor{qaline}{HTML}{2B303D}    % hairline rules / table rules
\definecolor{qatext}{HTML}{ECEEF2}    % primary text (off-white)
\definecolor{qamuted}{HTML}{8B93A6}   % secondary / caption text
\definecolor{qaaccent}{HTML}{D93ED7}  % primary accent -- violet du logo IUT, eclairci pour fond sombre
\definecolor{qaaccent2}{HTML}{F2B84B} % secondary accent -- or/ambre
\definecolor{qacodekw}{HTML}{C792EA}  % code : mots-cles
\definecolor{qacodestr}{HTML}{C3E88D} % code : chaines
\definecolor{qacodecom}{HTML}{7F8AA3} % code : commentaires
\definecolor{qacodeann}{HTML}{F2B84B} % code : annotations / javadoc

\setbeamercolor{background canvas}{bg=qabg}
\setbeamercolor{normal text}{fg=qatext}
\setbeamercolor{palette primary}{bg=qabg, fg=qamuted}
\setbeamercolor{frametitle}{bg=qabg, fg=qatext}
\setbeamercolor{progress bar}{fg=qaaccent, bg=qaline}
\setbeamercolor{alerted text}{fg=qaaccent2}
\setbeamercolor{block title}{bg=qapanel, fg=qaaccent}
\setbeamercolor{block body}{bg=qapanel, fg=qatext}
\setbeamercolor{block title alerted}{bg=qapanel, fg=qaaccent2}
\setbeamercolor{block body alerted}{bg=qapanel, fg=qatext}
\setbeamercolor{block title example}{bg=qapanelalt, fg=qaaccent}
\setbeamercolor{block body example}{bg=qapanelalt, fg=qatext}
\setbeamercolor{example text}{fg=qaaccent}
\setbeamercolor{itemize item}{fg=qaaccent}
\setbeamercolor{itemize subitem}{fg=qaaccent2}
\setbeamercolor{enumerate item}{fg=qaaccent}
\setbeamercolor{enumerate subitem}{fg=qaaccent2}
\setbeamercolor{caption}{fg=qamuted}
\setbeamercolor{page number in head/foot}{fg=qamuted, bg=qabg}
\setbeamercolor{title}{fg=qatext}
\setbeamercolor{subtitle}{fg=qamuted}
\setbeamercolor{author}{fg=qatext}
\setbeamercolor{institute}{fg=qamuted}
\setbeamercolor{date}{fg=qamuted}
\setbeamercolor{section in toc}{fg=qaaccent}
\setbeamercolor{subsection in toc}{fg=qamuted}
% Sommaire de debut de section : section en cours en violet et gras, les autres en gris
\setbeamercolor{section in toc shaded}{fg=qamuted}
\setbeamercolor{subsection in toc shaded}{fg=qamuted}

\setbeamertemplate{blocks}[rounded][shadow=false]
\setbeamertemplate{itemize item}{\raisebox{0.15em}{\scriptsize\textcolor{qaaccent}{\ensuremath{\blacktriangleright}}}}
\setbeamertemplate{itemize subitem}{\raisebox{0.1em}{\tiny\textcolor{qaaccent2}{\ensuremath{\bullet}}}}

\setbeamerfont{frametitle}{series=\bfseries, size=\Large}
\setbeamerfont{title}{series=\bfseries, size=\Large}
\setbeamerfont{subtitle}{size=\normalsize}
\setbeamerfont{institute}{size=\scriptsize}
\setbeamerfont{date}{size=\scriptsize}
\setbeamerfont{author}{size=\small}

\metroset{sectionpage=none, subsectionpage=none, numbering=fraction, progressbar=frametitle}
\setbeamertemplate{section in toc}[sections numbered]

% ---------------------------------------------------------------
% LISTINGS -- theme sombre (lstlisting reste utilisable tel quel)
% \lstcode : police + couleur de base, a reutiliser dans les
%            surcharges locales (ex. basicstyle=\lstcode\scriptsize)
% ---------------------------------------------------------------
\newcommand{\lstcode}{\ttfamily\color{qatext}}

\lstdefinelanguage{TypeScript}{
  keywords={typeof, new, true, false, catch, function, return, null, switch, var,
    if, in, while, do, else, case, break, const, let, interface, class, extends,
    implements, enum, export, import, this, super, constructor, public, private,
    protected, static, async, await, throw, throws, try, finally, void},
  ndkeywords={boolean, any, number, string, int, long, double},
  sensitive=true,
  comment=[l]{//},
  morecomment=[s]{/*}{*/},
  morestring=[b]',
  morestring=[b]"
}

\lstset{
  language=Java,
  basicstyle=\lstcode\footnotesize,
  keywordstyle=\color{qacodekw}\bfseries,
  ndkeywordstyle=\color{qaaccent2}\bfseries,
  identifierstyle=\color{qatext},
  stringstyle=\color{qacodestr},
  commentstyle=\color{qacodecom}\itshape,
  morecomment=[s][\color{qacodeann}]{/**}{*/},
  backgroundcolor=\color{qapanelalt},
  rulecolor=\color{qaline},
  frame=single,
  framesep=3pt,
  xleftmargin=4pt,
  xrightmargin=4pt,
  keepspaces=true,
  columns=fullflexible,
  showstringspaces=false,
  breaklines=true,
  tabsize=2,
  escapeinside={(*@}{@*)},
  % accents UTF-8 dans le code (listings ne gere pas l'UTF-8 multi-octets)
  literate={é}{{\'e}}1 {è}{{\`e}}1 {ê}{{\^e}}1 {à}{{\`a}}1 {â}{{\^a}}1
           {ù}{{\`u}}1 {û}{{\^u}}1 {ô}{{\^o}}1 {î}{{\^i}}1 {ç}{{\c c}}1
           {É}{{\'E}}1 {À}{{\`A}}1 {€}{{\euro{}}}1 {…}{{\ldots}}1,
}

% ---------------------------------------------------------------
% CARD-STYLE BOXES
% ---------------------------------------------------------------
\tcbset{
  figcard/.style={colback=white, colframe=white, boxrule=0pt, arc=4pt,
    left=3pt, right=3pt, top=3pt, bottom=3pt, boxsep=0pt,
    halign=center, valign=center},
  datacard/.style={colback=qapanel, colframe=qaline, boxrule=0.4pt, arc=3pt,
    left=8pt, right=8pt, top=6pt, bottom=6pt, coltext=qatext,
    before upper={\color{qatext}}, halign=center},
  codebox/.style={colback=qapanelalt, colframe=qaline, boxrule=0.4pt, arc=3pt,
    left=8pt, right=8pt, top=5pt, bottom=5pt, listing only,
    listing options={basicstyle=\ttfamily\small\color{qaaccent},
      breaklines=true, showstringspaces=false, columns=fullflexible,
      keepspaces=true, frame=none, backgroundcolor={}}}
}
\newtcolorbox{figcard}[1][]{figcard,#1}
\newtcolorbox{datacard}{datacard}
\newtcblisting{codebox}{codebox}

% ---------------------------------------------------------------
% HANDY MACROS
% ---------------------------------------------------------------
% Mention d'elaboration (discrete : plus sombre que le texte secondaire)
\definecolor{qaaimention}{HTML}{4E5566}
\newcommand{\raimention}{{\tiny\color{qaaimention}Supports \'elabor\'es avec l'aide de Claude Code Opus 5.5}}
\newcommand{\hl}[1]{\alert{#1}}
\newcommand{\source}[1]{{\scriptsize\color{qamuted}#1}}

% Sommaire automatique en debut de section
\AtBeginSection[]{%
  \begin{frame}<beamer>{Plan du cours}
    \tableofcontents[currentsection]
  \end{frame}}

% Page de titre : \rtitleslide{numero du cours}{titre du cours}
% Les auteurs sont repris de \rauthors (modifiable avant \begin{document})
\newcommand{\rauthors}{Chouki Tibermacine}
\newcommand{\rtitleslide}[2]{%
\begin{frame}[plain,noframenumbering]
  \begin{center}
    \vspace{2mm}
    {\color{qaaccent}\rule{2.1cm}{1.2pt}}\\[6mm]
    {\Large\bfseries\color{qatext} R4.02}\\[1mm]
    {\Large\bfseries\color{qatext} Qualit\'e de d\'eveloppement}\\[5mm]
    {\large\color{qamuted} Cours #1}\\[1mm]
    {\large\color{qaaccent} #2}\\[6mm]
    \includegraphics[height=2.1cm]{iut_logo_transparent.png}

    \vspace{4mm}

    {\small\color{qatext} \rauthors}\\[1mm]
    {\scriptsize\color{qamuted} BUT 2 Informatique --- IUT de Vannes --- Universit\'e Bretagne Sud}\\[2mm]
    \raimention
  \end{center}
\end{frame}
}

% ---- Diapo de demonstration : \rdemo{numero}{titre}{contenu (items)}{fichier(s) a ouvrir} ----
\newcommand{\rdemo}[4]{%
\begin{frame}{D\'emo #1 --- #2}
  \begin{tikzpicture}[remember picture, overlay]
    \node[anchor=north east, fill=qaaccent2, text=qabg, font=\scriptsize\bfseries,
      rounded corners=3pt, inner sep=4pt] at ([xshift=-0.4cm, yshift=-0.35cm]current page.north east) {D\'EMO};
  \end{tikzpicture}
  \begin{columns}[T]
    \begin{column}{0.62\textwidth}
      \begin{itemize}
        #3
      \end{itemize}
    \end{column}
    \begin{column}{0.34\textwidth}
      \begin{datacard}
        \raggedright\scriptsize
        {\color{qamuted} Dans l'IDE}\\[1mm]
        #4
      \end{datacard}
    \end{column}
  \end{columns}
\end{frame}}

% ---- Diapo de fin « Questions ? » ----
\newcommand{\rcontact}{chouki.tibermacine@univ-ubs.fr}
\newcommand{\rquestionslide}{%
\begin{frame}{Questions ?}
  \begin{center}
    \vspace{2mm}
    \begin{tikzpicture}
      \node[regular polygon, regular polygon sides=6, draw=qaline, thick,
        minimum size=4.6cm, rotate=90] (hexin) {};
      \node[regular polygon, regular polygon sides=6, draw=qaaccent, dashed,
        thick, minimum size=6.2cm, rotate=90] (hexout) {};
      \node[font=\fontsize{92}{92}\selectfont\bfseries, color=qaaccent2] at (hexin) {?};
    \end{tikzpicture}\\[5mm]
    {\small\color{qamuted} \rcontact}
  \end{center}
\end{frame}}

%   \graphicspath{{../../latex/common/assets/}{img/}{figs/}}
% ---------------------------------------------------------------
\usetheme{metropolis}
\usepackage{appendixnumberbeamer}

\usepackage[T1]{fontenc}
\usepackage[utf8]{inputenc}
\usepackage[french]{babel}
\usepackage[scale=0.94]{tgheros}
\renewcommand{\familydefault}{\sfdefault}
\usepackage{amsmath,amssymb,amsfonts}
\usepackage{graphicx}
\usepackage{booktabs}
\usepackage{multirow}
\usepackage{tikz}
\usetikzlibrary{shapes.geometric, shapes, arrows, arrows.meta, positioning, calc, fit}
\usepackage{xcolor}
\usepackage{listings}
\usepackage[skins,listings]{tcolorbox}
\usepackage{qrcode}
\usepackage{eurosym}
\usepackage{hyperref}

% ---------------------------------------------------------------
% PALETTE (dark-mode tech: near-black canvas, violet (logo IUT) + gold accents)
% ---------------------------------------------------------------
\definecolor{qabg}{HTML}{0E1016}      % canvas background
\definecolor{qapanel}{HTML}{181B24}   % card / block panel
\definecolor{qapanelalt}{HTML}{20242F}% secondary panel tint
\definecolor{qaline}{HTML}{2B303D}    % hairline rules / table rules
\definecolor{qatext}{HTML}{ECEEF2}    % primary text (off-white)
\definecolor{qamuted}{HTML}{8B93A6}   % secondary / caption text
\definecolor{qaaccent}{HTML}{D93ED7}  % primary accent -- violet du logo IUT, eclairci pour fond sombre
\definecolor{qaaccent2}{HTML}{F2B84B} % secondary accent -- or/ambre
\definecolor{qacodekw}{HTML}{C792EA}  % code : mots-cles
\definecolor{qacodestr}{HTML}{C3E88D} % code : chaines
\definecolor{qacodecom}{HTML}{7F8AA3} % code : commentaires
\definecolor{qacodeann}{HTML}{F2B84B} % code : annotations / javadoc

\setbeamercolor{background canvas}{bg=qabg}
\setbeamercolor{normal text}{fg=qatext}
\setbeamercolor{palette primary}{bg=qabg, fg=qamuted}
\setbeamercolor{frametitle}{bg=qabg, fg=qatext}
\setbeamercolor{progress bar}{fg=qaaccent, bg=qaline}
\setbeamercolor{alerted text}{fg=qaaccent2}
\setbeamercolor{block title}{bg=qapanel, fg=qaaccent}
\setbeamercolor{block body}{bg=qapanel, fg=qatext}
\setbeamercolor{block title alerted}{bg=qapanel, fg=qaaccent2}
\setbeamercolor{block body alerted}{bg=qapanel, fg=qatext}
\setbeamercolor{block title example}{bg=qapanelalt, fg=qaaccent}
\setbeamercolor{block body example}{bg=qapanelalt, fg=qatext}
\setbeamercolor{example text}{fg=qaaccent}
\setbeamercolor{itemize item}{fg=qaaccent}
\setbeamercolor{itemize subitem}{fg=qaaccent2}
\setbeamercolor{enumerate item}{fg=qaaccent}
\setbeamercolor{enumerate subitem}{fg=qaaccent2}
\setbeamercolor{caption}{fg=qamuted}
\setbeamercolor{page number in head/foot}{fg=qamuted, bg=qabg}
\setbeamercolor{title}{fg=qatext}
\setbeamercolor{subtitle}{fg=qamuted}
\setbeamercolor{author}{fg=qatext}
\setbeamercolor{institute}{fg=qamuted}
\setbeamercolor{date}{fg=qamuted}
\setbeamercolor{section in toc}{fg=qaaccent}
\setbeamercolor{subsection in toc}{fg=qamuted}
% Sommaire de debut de section : section en cours en violet et gras, les autres en gris
\setbeamercolor{section in toc shaded}{fg=qamuted}
\setbeamercolor{subsection in toc shaded}{fg=qamuted}

\setbeamertemplate{blocks}[rounded][shadow=false]
\setbeamertemplate{itemize item}{\raisebox{0.15em}{\scriptsize\textcolor{qaaccent}{\ensuremath{\blacktriangleright}}}}
\setbeamertemplate{itemize subitem}{\raisebox{0.1em}{\tiny\textcolor{qaaccent2}{\ensuremath{\bullet}}}}

\setbeamerfont{frametitle}{series=\bfseries, size=\Large}
\setbeamerfont{title}{series=\bfseries, size=\Large}
\setbeamerfont{subtitle}{size=\normalsize}
\setbeamerfont{institute}{size=\scriptsize}
\setbeamerfont{date}{size=\scriptsize}
\setbeamerfont{author}{size=\small}

\metroset{sectionpage=none, subsectionpage=none, numbering=fraction, progressbar=frametitle}
\setbeamertemplate{section in toc}[sections numbered]

% ---------------------------------------------------------------
% LISTINGS -- theme sombre (lstlisting reste utilisable tel quel)
% \lstcode : police + couleur de base, a reutiliser dans les
%            surcharges locales (ex. basicstyle=\lstcode\scriptsize)
% ---------------------------------------------------------------
\newcommand{\lstcode}{\ttfamily\color{qatext}}

\lstdefinelanguage{TypeScript}{
  keywords={typeof, new, true, false, catch, function, return, null, switch, var,
    if, in, while, do, else, case, break, const, let, interface, class, extends,
    implements, enum, export, import, this, super, constructor, public, private,
    protected, static, async, await, throw, throws, try, finally, void},
  ndkeywords={boolean, any, number, string, int, long, double},
  sensitive=true,
  comment=[l]{//},
  morecomment=[s]{/*}{*/},
  morestring=[b]',
  morestring=[b]"
}

\lstset{
  language=Java,
  basicstyle=\lstcode\footnotesize,
  keywordstyle=\color{qacodekw}\bfseries,
  ndkeywordstyle=\color{qaaccent2}\bfseries,
  identifierstyle=\color{qatext},
  stringstyle=\color{qacodestr},
  commentstyle=\color{qacodecom}\itshape,
  morecomment=[s][\color{qacodeann}]{/**}{*/},
  backgroundcolor=\color{qapanelalt},
  rulecolor=\color{qaline},
  frame=single,
  framesep=3pt,
  xleftmargin=4pt,
  xrightmargin=4pt,
  keepspaces=true,
  columns=fullflexible,
  showstringspaces=false,
  breaklines=true,
  tabsize=2,
  escapeinside={(*@}{@*)},
  % accents UTF-8 dans le code (listings ne gere pas l'UTF-8 multi-octets)
  literate={é}{{\'e}}1 {è}{{\`e}}1 {ê}{{\^e}}1 {à}{{\`a}}1 {â}{{\^a}}1
           {ù}{{\`u}}1 {û}{{\^u}}1 {ô}{{\^o}}1 {î}{{\^i}}1 {ç}{{\c c}}1
           {É}{{\'E}}1 {À}{{\`A}}1 {€}{{\euro{}}}1 {…}{{\ldots}}1,
}

% ---------------------------------------------------------------
% CARD-STYLE BOXES
% ---------------------------------------------------------------
\tcbset{
  figcard/.style={colback=white, colframe=white, boxrule=0pt, arc=4pt,
    left=3pt, right=3pt, top=3pt, bottom=3pt, boxsep=0pt,
    halign=center, valign=center},
  datacard/.style={colback=qapanel, colframe=qaline, boxrule=0.4pt, arc=3pt,
    left=8pt, right=8pt, top=6pt, bottom=6pt, coltext=qatext,
    before upper={\color{qatext}}, halign=center},
  codebox/.style={colback=qapanelalt, colframe=qaline, boxrule=0.4pt, arc=3pt,
    left=8pt, right=8pt, top=5pt, bottom=5pt, listing only,
    listing options={basicstyle=\ttfamily\small\color{qaaccent},
      breaklines=true, showstringspaces=false, columns=fullflexible,
      keepspaces=true, frame=none, backgroundcolor={}}}
}
\newtcolorbox{figcard}[1][]{figcard,#1}
\newtcolorbox{datacard}{datacard}
\newtcblisting{codebox}{codebox}

% ---------------------------------------------------------------
% HANDY MACROS
% ---------------------------------------------------------------
% Mention d'elaboration (discrete : plus sombre que le texte secondaire)
\definecolor{qaaimention}{HTML}{4E5566}
\newcommand{\raimention}{{\tiny\color{qaaimention}Supports \'elabor\'es avec l'aide de Claude Code Opus 5.5}}
\newcommand{\hl}[1]{\alert{#1}}
\newcommand{\source}[1]{{\scriptsize\color{qamuted}#1}}

% Sommaire automatique en debut de section
\AtBeginSection[]{%
  \begin{frame}<beamer>{Plan du cours}
    \tableofcontents[currentsection]
  \end{frame}}

% Page de titre : \rtitleslide{numero du cours}{titre du cours}
% Les auteurs sont repris de \rauthors (modifiable avant \begin{document})
\newcommand{\rauthors}{Chouki Tibermacine}
\newcommand{\rtitleslide}[2]{%
\begin{frame}[plain,noframenumbering]
  \begin{center}
    \vspace{2mm}
    {\color{qaaccent}\rule{2.1cm}{1.2pt}}\\[6mm]
    {\Large\bfseries\color{qatext} R4.02}\\[1mm]
    {\Large\bfseries\color{qatext} Qualit\'e de d\'eveloppement}\\[5mm]
    {\large\color{qamuted} Cours #1}\\[1mm]
    {\large\color{qaaccent} #2}\\[6mm]
    \includegraphics[height=2.1cm]{iut_logo_transparent.png}

    \vspace{4mm}

    {\small\color{qatext} \rauthors}\\[1mm]
    {\scriptsize\color{qamuted} BUT 2 Informatique --- IUT de Vannes --- Universit\'e Bretagne Sud}\\[2mm]
    \raimention
  \end{center}
\end{frame}
}

% ---- Diapo de demonstration : \rdemo{numero}{titre}{contenu (items)}{fichier(s) a ouvrir} ----
\newcommand{\rdemo}[4]{%
\begin{frame}{D\'emo #1 --- #2}
  \begin{tikzpicture}[remember picture, overlay]
    \node[anchor=north east, fill=qaaccent2, text=qabg, font=\scriptsize\bfseries,
      rounded corners=3pt, inner sep=4pt] at ([xshift=-0.4cm, yshift=-0.35cm]current page.north east) {D\'EMO};
  \end{tikzpicture}
  \begin{columns}[T]
    \begin{column}{0.62\textwidth}
      \begin{itemize}
        #3
      \end{itemize}
    \end{column}
    \begin{column}{0.34\textwidth}
      \begin{datacard}
        \raggedright\scriptsize
        {\color{qamuted} Dans l'IDE}\\[1mm]
        #4
      \end{datacard}
    \end{column}
  \end{columns}
\end{frame}}

% ---- Diapo de fin « Questions ? » ----
\newcommand{\rcontact}{chouki.tibermacine@univ-ubs.fr}
\newcommand{\rquestionslide}{%
\begin{frame}{Questions ?}
  \begin{center}
    \vspace{2mm}
    \begin{tikzpicture}
      \node[regular polygon, regular polygon sides=6, draw=qaline, thick,
        minimum size=4.6cm, rotate=90] (hexin) {};
      \node[regular polygon, regular polygon sides=6, draw=qaaccent, dashed,
        thick, minimum size=6.2cm, rotate=90] (hexout) {};
      \node[font=\fontsize{92}{92}\selectfont\bfseries, color=qaaccent2] at (hexin) {?};
    \end{tikzpicture}\\[5mm]
    {\small\color{qamuted} \rcontact}
  \end{center}
\end{frame}}

\graphicspath{{../../latex/common/assets/}{img/}}

\title{Qualité logicielle et panorama des tests}
\author[C. Tibermacine]{Chouki Tibermacine}
\date{}

\lstdefinestyle{sortie}{language={}, basicstyle=\lstcode\scriptsize, keywordstyle={}, commentstyle={}, stringstyle={},
  literate={é}{{\'e}}1 {è}{{\`e}}1 {ô}{{\^o}}1 {€}{{\euro{}}}1
           {✔}{{\textcolor{qacodestr}{$\checkmark$}}}1 {✘}{{\textcolor{qaaccent}{$\times$}}}1}

\tikzset{
  layer/.style={rectangle, rounded corners=2pt, draw=qaaccent, thick, fill=qapanel, text=qatext,
    minimum height=0.55cm, minimum width=3.6cm, align=center, font=\scriptsize},
  ext/.style={rectangle, rounded corners=2pt, draw=qaaccent2, thick, dashed, fill=qapanelalt, text=qatext,
    minimum height=0.5cm, minimum width=2.4cm, align=center, font=\scriptsize},
  tag/.style={rectangle, rounded corners=6pt, fill=qaaccent2, text=qabg, font=\tiny\bfseries,
    inner sep=2pt},
  arr/.style={-Latex, qamuted, thick}
}

\begin{document}

\rtitleslide{1}{Qualité logicielle et panorama des tests}

\begin{frame}{Objectifs du module R4.02}
  \begin{itemize}
    \item Aller plus loin que les tests unitaires JUnit de BUT1 :
      \begin{itemize}
        \item partir du \textbf{besoin} (user stories) pour écrire des \textbf{tests d'acceptation}
        \item \textbf{isoler} le code testé avec des doublures (\textit{mocks})
        \item tester l'\textbf{intégration} avec une vraie base de données
        \item tester l'application \textbf{de bout en bout}, depuis le navigateur
        \item faire exécuter tous ces tests \textbf{automatiquement} à chaque modification (CI)
      \end{itemize}
    \item Savoir juger si un projet (ou une fonctionnalité) est \textbf{faisable} et \textbf{testable}
    \item Un fil rouge : l'application bancaire que vous connaissez depuis R3.07
  \end{itemize}
  \vspace{3mm}
  \begin{datacard}
    \small user story $\rightarrow$ critères d'acceptation $\rightarrow$ tests unitaires $\rightarrow$ tests d'intégration $\rightarrow$ tests de bout en bout $\rightarrow$ pipeline de CI
  \end{datacard}
\end{frame}

% ================================================================
\section{Pourquoi tester ?}

\begin{frame}{Trois incidents célèbres}
  \begin{columns}[T]
    \begin{column}{0.32\textwidth}
      \begin{block}{Ariane 5, vol 501 (1996)}
        \footnotesize
        \begin{itemize}
          \item Explosion 37 s après le décollage
          \item Un nombre flottant 64 bits converti en entier 16 bits : \alert{dépassement}
          \item Code \textbf{réutilisé d'Ariane 4}, jamais testé avec la trajectoire d'Ariane 5
        \end{itemize}
      \end{block}
    \end{column}
    \begin{column}{0.32\textwidth}
      \begin{block}{Knight Capital (2012)}
        \footnotesize
        \begin{itemize}
          \item Société de trading : \alert{440 M\$ perdus en 45 min}
          \item Nouveau code déployé sur 7 serveurs sur 8 ; le 8\textsuperscript{e} exécute un \textbf{vieux code mort} réactivé par un drapeau réutilisé
          \item Déploiement manuel, sans vérification automatique
        \end{itemize}
      \end{block}
    \end{column}
    \begin{column}{0.32\textwidth}
      \begin{block}{CrowdStrike (juillet 2024)}
        \footnotesize
        \begin{itemize}
          \item Une mise à jour de configuration d'un logiciel de sécurité
          \item \alert{8,5 millions de PC Windows} bloqués (écran bleu) : aéroports, hôpitaux, banques
          \item Lecture hors des limites d'un tableau, non détectée par le validateur ; mise à jour envoyée à \textbf{tout le monde en même temps}
        \end{itemize}
      \end{block}
    \end{column}
  \end{columns}
  \vspace{3mm}
  \begin{center}
    \small Points communs : du code \textbf{réutilisé} ou \textbf{modifié} sans être re-testé dans son nouveau contexte, et des déploiements sans filet.
  \end{center}
\end{frame}

\begin{frame}{Plus un bug est trouvé tard, plus il coûte cher}
  \begin{columns}[c]
    \begin{column}{0.58\textwidth}
      \begin{center}
      \begin{tikzpicture}[x=1cm, y=1cm]
        \draw[qaline] (0,0) -- (6.6,0);
        \foreach \i/\h/\l in {0/0.25/Exigences, 1/0.45/Conception, 2/0.75/Code, 3/1.3/Tests, 4/2.1/Recette, 5/3.4/Production} {
          \fill[qaaccent!70!qapanel] (\i*1.1+0.15,0) rectangle (\i*1.1+0.85,\h);
          \node[font=\tiny, text=qamuted, anchor=north] at (\i*1.1+0.5,-0.05) {\l};
        }
        \fill[qaaccent2] (5*1.1+0.15,0) rectangle (5*1.1+0.85,3.4);
        \node[font=\scriptsize, text=qamuted, rotate=90, anchor=south] at (-0.15,1.7) {coût de correction};
      \end{tikzpicture}
      \end{center}
      \source{Ordre de grandeur popularisé par Boehm (1981). Les chiffres exacts sont discutés, la tendance ne l'est pas.}
    \end{column}
    \begin{column}{0.40\textwidth}
      \begin{itemize}
        \item En production, un bug coûte :
          \begin{itemize}
            \item le diagnostic (souvent sans accès au contexte)
            \item la correction \textbf{et} un nouveau déploiement
            \item les données à réparer, la confiance des clients
          \end{itemize}
        \item Tester tôt et souvent $=$ \alert{décaler la détection vers la gauche} (\textit{shift left})
      \end{itemize}
    \end{column}
  \end{columns}
\end{frame}

\begin{frame}{Qu'est-ce qu'un logiciel de qualité ? (ISO/IEC 25010)}
  \small
  \begin{columns}[T]
    \begin{column}{0.32\textwidth}
      \begin{datacard}
        \textbf{\color{qaaccent}Adéquation fonctionnelle}\\[-0.5mm]
        {\scriptsize\color{qamuted} fait ce qui est demandé, correctement}
      \end{datacard}
      \begin{datacard}
        \textbf{\color{qaaccent}Fiabilité}\\[-0.5mm]
        {\scriptsize\color{qamuted} fonctionne dans la durée, résiste aux pannes}
      \end{datacard}
      \begin{datacard}
        \textbf{Efficacité de performance}\\[-0.5mm]
        {\scriptsize\color{qamuted} temps de réponse, ressources}
      \end{datacard}
    \end{column}
    \begin{column}{0.32\textwidth}
      \begin{datacard}
        \textbf{\color{qaaccent}Maintenabilité}\\[-0.5mm]
        {\scriptsize\color{qamuted} facile à modifier\ldots{} et à \textbf{tester}}
      \end{datacard}
      \begin{datacard}
        \textbf{Sécurité}\\[-0.5mm]
        {\scriptsize\color{qamuted} confidentialité, intégrité}
      \end{datacard}
      \begin{datacard}
        \textbf{Capacité d'interaction}\\[-0.5mm]
        {\scriptsize\color{qamuted} utilisabilité, accessibilité}
      \end{datacard}
    \end{column}
    \begin{column}{0.32\textwidth}
      \begin{datacard}
        \textbf{Compatibilité}\\[-0.5mm]
        {\scriptsize\color{qamuted} coexiste et échange avec d'autres systèmes}
      \end{datacard}
      \begin{datacard}
        \textbf{Flexibilité}\\[-0.5mm]
        {\scriptsize\color{qamuted} s'adapte, s'installe ailleurs}
      \end{datacard}
      \begin{datacard}
        \textbf{Sûreté}\\[-0.5mm]
        {\scriptsize\color{qamuted} ne met personne en danger}
      \end{datacard}
    \end{column}
  \end{columns}
  \vspace{2mm}
  \begin{center}
    Les tests touchent presque toutes ces caractéristiques ; ce module insiste sur celles \textcolor{qaaccent}{en violet}.
  \end{center}
\end{frame}

\begin{frame}{Vérification et validation}
  \begin{columns}[T]
    \begin{column}{0.48\textwidth}
      \begin{block}{Vérification}
        \textit{« Construisons-nous le produit correctement ? »}
        \begin{itemize}
          \item Le code respecte-t-il sa spécification ?
          \item Tests unitaires, tests d'intégration
          \item Question posée aux \textbf{développeurs}
        \end{itemize}
      \end{block}
    \end{column}
    \begin{column}{0.48\textwidth}
      \begin{block}{Validation}
        \textit{« Construisons-nous le bon produit ? »}
        \begin{itemize}
          \item Le produit répond-il au besoin réel ?
          \item Tests d'acceptation, recette, tests utilisateurs
          \item Question posée au \textbf{client}
        \end{itemize}
      \end{block}
    \end{column}
  \end{columns}
  \vspace{4mm}
  \begin{itemize}
    \item Un code parfaitement vérifié peut être inutile s'il implémente le mauvais besoin
    \item Le CM2 partira du besoin (user stories) pour écrire des tests de \textbf{validation}
  \end{itemize}
\end{frame}

\begin{frame}{Ce qu'un test peut prouver\ldots{} et ce qu'il ne peut pas}
  \begin{center}
    \begin{datacard}
      \textit{« Tester un programme peut montrer la présence de bugs,\\ jamais leur absence. »}\\[1mm]
      {\scriptsize\color{qamuted} Edsger W. Dijkstra, 1970}
    \end{datacard}
  \end{center}
  \begin{itemize}
    \item Tester tous les cas est impossible : une méthode \texttt{f(int a, int b)} a $2^{64} \approx 1{,}8 \times 10^{19}$ combinaisons d'entrées
    \item Tester, c'est donc \textbf{choisir} un petit nombre de cas représentatifs
      \begin{itemize}
        \item les cas \textbf{nominaux} (ce qui doit marcher)
        \item les cas \textbf{limites} (0, solde exact, liste vide\ldots)
        \item les cas d'\textbf{erreur} (montant négatif, compte inexistant\ldots)
      \end{itemize}
    \item Un test vert donne de la \textbf{confiance}, pas une preuve
  \end{itemize}
\end{frame}

% ================================================================
\section{Panorama des tests}

\begin{frame}{Les niveaux de tests}
  \small
  \begin{datacard}
    \begin{tabular}{@{}p{2.6cm}p{6.0cm}p{3.6cm}@{}}
      \toprule
      \textbf{Niveau} & \textbf{Ce qu'on teste} & \textbf{Dans ce module} \\
      \midrule
      \textbf{Unitaire} & une classe ou une méthode, \textbf{isolée} du reste & JUnit, Mockito (TP1, TP3) \\
      \textbf{Intégration} & plusieurs composants ensemble : code + base de données, code + API & Testcontainers, Spring (TP4) \\
      \textbf{Système / bout en bout} & l'application complète, comme un utilisateur, depuis le navigateur & Playwright (TP5) \\
      \textbf{Acceptation} & le besoin du client est-il satisfait ? & Cucumber (TP2), recette \\
      \bottomrule
    \end{tabular}
  \end{datacard}
  \vspace{2mm}
  \begin{itemize}
    \item Plus on monte : plus c'est \textbf{réaliste}, mais plus c'est \textbf{lent}, \textbf{fragile} et \textbf{difficile à diagnostiquer}
  \end{itemize}
\end{frame}

\begin{frame}{Les types de tests}
  \begin{columns}[T]
    \begin{column}{0.48\textwidth}
      \begin{block}{Selon ce qu'on vérifie}
        \begin{itemize}
          \item \textbf{Fonctionnels} : le résultat est-il correct ?
          \item \textbf{Non fonctionnels} :
            \begin{itemize}
              \item performance, montée en charge
              \item sécurité
              \item utilisabilité, accessibilité
            \end{itemize}
        \end{itemize}
      \end{block}
    \end{column}
    \begin{column}{0.48\textwidth}
      \begin{block}{Selon ce qu'on connaît du code}
        \begin{itemize}
          \item \textbf{Boîte noire} : on ne regarde que la spécification (entrées et sorties attendues)
          \item \textbf{Boîte blanche} : on s'appuie sur le code (chaque branche d'un \texttt{if} est-elle exécutée ?)
            \begin{itemize}
              \item mesuré par la \textbf{couverture} (CM6)
            \end{itemize}
        \end{itemize}
      \end{block}
    \end{column}
  \end{columns}
  \vspace{3mm}
  \begin{itemize}
    \item Les deux axes se combinent : un test de bout en bout est en général fonctionnel et boîte noire
  \end{itemize}
\end{frame}

\begin{frame}{Quel niveau ? Quel type ?}
  \small
  \begin{enumerate}
    \item Vérifier que \texttt{calculerInterets()} renvoie 2\,\% du solde
      \uncover<2->{\hfill \alert{unitaire, fonctionnel}}
    \item Vérifier que \texttt{CompteRepository.findByClient()} renvoie les bons comptes depuis PostgreSQL
      \uncover<2->{\hfill \alert{intégration}}
    \item Se connecter dans le navigateur, faire un virement, voir le nouveau solde affiché
      \uncover<2->{\hfill \alert{bout en bout}}
    \item Le conseiller vérifie la règle « un virement de plus de 1\,000\,\euro{} demande une confirmation »
      \uncover<2->{\hfill \alert{acceptation}}
    \item La page des comptes s'affiche en moins de 2 s avec 1\,000 utilisateurs connectés
      \uncover<2->{\hfill \alert{non fonctionnel (performance)}}
    \item Après une correction dans \texttt{Compte}, relancer tous les tests existants
      \uncover<2->{\hfill \alert{non-régression}}
  \end{enumerate}
\end{frame}

\begin{frame}{La pyramide des tests}
  \begin{columns}[c]
    \begin{column}{0.5\textwidth}
      \begin{center}
      \begin{tikzpicture}[x=0.85cm, y=1cm]
        \fill[qaaccent!35!qapanel] (0,0) -- (6,0) -- (4.6,1.4) -- (1.4,1.4) -- cycle;
        \fill[qaaccent!65!qapanel] (1.4,1.4) -- (4.6,1.4) -- (3.7,2.3) -- (2.3,2.3) -- cycle;
        \fill[qaaccent2] (2.3,2.3) -- (3.7,2.3) -- (3,3.0) -- cycle;
        \node[font=\small\bfseries, text=qatext] at (3,0.7) {Unitaires};
        \node[font=\scriptsize\bfseries, text=qatext] at (3,1.85) {Intégration};
        \node[font=\tiny\bfseries, text=qabg] at (3,2.55) {E2E};
        \draw[arr] (6.4,0.1) -- node[right, font=\tiny, text=qamuted, align=left] {plus lents\\plus coûteux\\plus fragiles} (6.4,2.9);
      \end{tikzpicture}
      \end{center}
      \source{Mike Cohn, \textit{Succeeding with Agile}, 2009}
    \end{column}
    \begin{column}{0.48\textwidth}
      \begin{itemize}
        \item \textbf{Beaucoup} de tests unitaires : rapides (des milliers en quelques secondes), précis quand ils échouent
        \item \textbf{Moins} de tests d'intégration : ils vérifient que les morceaux s'assemblent
        \item \textbf{Peu} de tests de bout en bout : les parcours essentiels de l'utilisateur
        \item Chaque bug doit être attrapé \alert{au niveau le plus bas possible}
      \end{itemize}
    \end{column}
  \end{columns}
\end{frame}

\begin{frame}{Variantes et anti-modèle}
  \begin{columns}[T]
    \begin{column}{0.48\textwidth}
      \begin{block}{Le trophée des tests}
        \begin{itemize}
          \item Proposé pour les applications web (Kent C. Dodds)
          \item Met l'accent sur les tests d'\textbf{intégration} : c'est souvent là que les bugs se cachent (base, API, configuration)
          \item Ajoute à la base l'\textbf{analyse statique} (typage, \textit{linter})
        \end{itemize}
      \end{block}
    \end{column}
    \begin{column}{0.48\textwidth}
      \begin{alertblock}{Le cornet de glace (anti-modèle)}
        \begin{itemize}
          \item La pyramide à l'envers : beaucoup de tests manuels et de bout en bout, presque pas de tests unitaires
          \item Suite lente (des heures), fragile, qu'on finit par ne plus lancer
          \item Fréquent quand on teste \og après coup \fg{} un code non conçu pour être testé
        \end{itemize}
      \end{alertblock}
    \end{column}
  \end{columns}
  \vspace{3mm}
  \begin{center}
    La bonne forme dépend de l'application ; le principe reste : \textbf{des retours rapides} sur la majorité des erreurs.
  \end{center}
\end{frame}

\begin{frame}{La non-régression}
  \small
  \begin{columns}[T]
    \begin{column}{0.52\textwidth}
      \begin{itemize}
        \item Une \textbf{régression} : une fonctionnalité qui marchait ne marche plus après une modification
        \item Exemple :
          \begin{enumerate}
            \item semaine 1 : on corrige le calcul du plafond de retrait
            \item semaine 3 : un collègue modifie le calcul des frais\ldots
            \item \ldots{} et le plafond est de nouveau faux
          \end{enumerate}
        \item Les \textbf{tests de non-régression} : rejouer \alert{tous} les tests existants après \alert{chaque} modification
      \end{itemize}
    \end{column}
    \begin{column}{0.44\textwidth}
      \begin{block}{Pourquoi automatiser ?}
        \begin{itemize}
          \item Rejouer à la main des centaines de cas à chaque modification : impossible
          \item Automatisés, les tests deviennent un \textbf{filet de sécurité} qui permet de modifier le code sans peur
          \item Lancés par la pipeline de CI à chaque \textit{push} (CM6)
        \end{itemize}
      \end{block}
      \begin{exampleblock}{Règle d'or}
        Un bug corrigé $=$ un test ajouté qui le reproduit
      \end{exampleblock}
    \end{column}
  \end{columns}
\end{frame}

% ================================================================
\section{JUnit : rappels et nouveautés}

\begin{frame}[fragile]{Ce que vous savez déjà}
  \begin{columns}[T]
    \begin{column}{0.55\textwidth}
\begin{lstlisting}[basicstyle=\lstcode\scriptsize]
class CompteTest {
    private Compte compte;

    @BeforeEach
    void setUp() {
        compte = new Compte("FR76-001", new BigDecimal("100.00"));
    }

    @Test
    void crediter_augmenteLeSolde() {
        compte.crediter(new BigDecimal("50.00"));
        assertEquals(new BigDecimal("150.00"), compte.getSolde());
    }
}
\end{lstlisting}
    \end{column}
    \begin{column}{0.42\textwidth}
      \begin{itemize}
        \item JUnit Jupiter (JUnit 5, puis JUnit 6) : \texttt{org.junit.jupiter.api}
        \item \texttt{@Test}, \texttt{@BeforeEach}, \texttt{@AfterEach}, \texttt{@BeforeAll}, \texttt{@AfterAll}
        \item \texttt{assertEquals}, \texttt{assertTrue}, \texttt{assertNull}\ldots
        \item Une \textbf{nouvelle instance} de la classe de test pour chaque méthode de test
        \item Lancement : IDE ou \texttt{mvn test}
      \end{itemize}
    \end{column}
  \end{columns}
\end{frame}

\begin{frame}[fragile]{Structurer un test : Arrange, Act, Assert}
  \begin{columns}[T]
    \begin{column}{0.6\textwidth}
\begin{lstlisting}[basicstyle=\lstcode\scriptsize]
@Test
void debiter_soldeSuffisant_diminueLeSolde() {
    // Arrange (Given) : préparer le contexte
    Compte compte = new Compte("FR76-001",
                               new BigDecimal("100.00"));

    // Act (When) : UNE action
    compte.debiter(new BigDecimal("30.00"));

    // Assert (Then) : vérifier le résultat
    assertEquals(new BigDecimal("70.00"),
                 compte.getSolde());
}
\end{lstlisting}
    \end{column}
    \begin{column}{0.37\textwidth}
      \begin{itemize}
        \item Trois blocs, toujours dans cet ordre
        \item Une seule action testée par test
        \item Équivalent BDD : \textbf{Given / When / Then} (CM2)
        \item Nommer le test pour qu'il se lise comme une phrase :
          \begin{itemize}
            \item {\scriptsize\texttt{methode\_contexte\_resultat}}
            \item ou \texttt{@DisplayName("...")}
          \end{itemize}
        \item Quand il échoue, le nom doit suffire à comprendre \textbf{ce qui} ne marche plus
      \end{itemize}
    \end{column}
  \end{columns}
\end{frame}

\begin{frame}{Les propriétés d'un bon test : FIRST}
  \begin{datacard}
    \begin{tabular}{@{}cl p{8.2cm}@{}}
      \textbf{\color{qaaccent}\Large F} & \textbf{Fast} & rapide : quelques millisecondes, pour pouvoir lancer toute la suite souvent \\[1mm]
      \textbf{\color{qaaccent}\Large I} & \textbf{Independent} & indépendant : ne dépend ni des autres tests, ni de leur ordre d'exécution \\[1mm]
      \textbf{\color{qaaccent}\Large R} & \textbf{Repeatable} & reproductible : même résultat partout et à chaque fois (pas de date du jour, de hasard, de réseau\ldots) \\[1mm]
      \textbf{\color{qaaccent}\Large S} & \textbf{Self-validating} & auto-vérifiant : vert ou rouge, sans lire de \texttt{println} \\[1mm]
      \textbf{\color{qaaccent}\Large T} & \textbf{Timely} & écrit au bon moment : avec le code (voire avant), pas des mois après \\
    \end{tabular}
  \end{datacard}
  \vspace{1mm}
  \begin{itemize}\small
    \item Les tests qui touchent une base, une horloge ou une API externe cassent souvent \textbf{I} et \textbf{R} : on verra comment les isoler (CM3) ou les maîtriser (CM4)
  \end{itemize}
\end{frame}

\begin{frame}{Choisir ses cas de test : partitions et valeurs limites}
  Spécification : \texttt{debiter(montant)} accepte un montant strictement positif et inférieur ou égal au solde.
  \vspace{2mm}
  \begin{center}
  \begin{tikzpicture}[x=1cm, y=1cm]
    \fill[qaaccent2!40!qapanel] (0,0) rectangle (3,0.6);
    \fill[qaaccent!45!qapanel] (3,0) rectangle (8,0.6);
    \fill[qaaccent2!40!qapanel] (8,0) rectangle (11,0.6);
    \node[font=\scriptsize, text=qatext] at (1.5,0.3) {montant $\leq 0$ : refusé};
    \node[font=\scriptsize, text=qatext] at (5.5,0.3) {$0 <$ montant $\leq$ solde : accepté};
    \node[font=\scriptsize, text=qatext] at (9.5,0.3) {montant $>$ solde : refusé};
    \foreach \x/\l in {2.35/-0.01, 3/0, 3.65/0.01, 7.35/99.99, 8/100.00, 8.65/100.01} {
      \draw[qaaccent2, thick] (\x,-0.05) -- (\x,0.65);
      \node[font=\tiny, text=qaaccent2, anchor=north] at (\x,-0.1) {\l};
    }
    \node[font=\tiny, text=qamuted] at (5.5,-0.55) {(solde = 100.00)};
  \end{tikzpicture}
  \end{center}
  \begin{itemize}
    \item \textbf{Partitions d'équivalence} : des classes d'entrées qui doivent se comporter pareil ; un cas par classe suffit
    \item \textbf{Valeurs limites} : les bugs se cachent aux frontières (\texttt{<} au lieu de \texttt{<=}) ; tester de chaque côté
    \item Ici : 3 partitions et 6 valeurs limites $\rightarrow$ un \textbf{test paramétré}
  \end{itemize}
\end{frame}

\begin{frame}[fragile]{Tests paramétrés}
\begin{lstlisting}[basicstyle=\lstcode\scriptsize]
@ParameterizedTest(name = "débit de {0} refusé")
@ValueSource(strings = {"0", "-0.01", "-10"})
void debiter_montantNonPositif_refuse(String montant) {
    Compte compte = new Compte("FR76-001", new BigDecimal("100.00"));
    assertThrows(IllegalArgumentException.class, () -> compte.debiter(new BigDecimal(montant)));
}

@ParameterizedTest(name = "{0} - {1} = {2}")
@CsvSource({
    "100.00,  0.01, 99.99",
    "100.00, 99.99,  0.01",
    "100.00, 100.00, 0.00"     // cas limite : le solde tombe à zéro
})
void debiter_soldeSuffisant(BigDecimal initial, BigDecimal montant, BigDecimal attendu) {
    Compte compte = new Compte("FR76-001", initial);
    compte.debiter(montant);
    assertEquals(attendu, compte.getSolde());
}
\end{lstlisting}
  \small Autres sources : \texttt{@EnumSource}, \texttt{@MethodSource} (objets complexes), \texttt{@CsvFileSource}
\end{frame}

\rdemo{1}{Une valeur limite attrape un bug}{%
  \item Lancer \texttt{CompteTest} : l'IDE affiche les tests comme une spécification
  \item On glisse un bug dans \texttt{debiter} : \texttt{<} devient \texttt{<=}
  \item Un seul test devient rouge : le cas \textbf{limite} où le solde tombe à 0
  \item Tous les cas « au milieu » restent verts : sans la valeur limite, le bug passait
}{\texttt{Compte}\\\texttt{CompteTest}}

\begin{frame}[fragile]{Tester les exceptions}
  \begin{columns}[T]
    \begin{column}{0.6\textwidth}
\begin{lstlisting}[basicstyle=\lstcode\scriptsize]
// eur("50.00") : raccourci pour new BigDecimal("50.00")

@Test
void debiter_soldeInsuffisant_refuse() {
  Compte compte = new Compte("FR76-001", eur("50.00"));

  SoldeInsuffisantException e = assertThrows(
      SoldeInsuffisantException.class,
      () -> compte.debiter(eur("80.00")));

  assertEquals("FR76-001", e.getNumeroCompte());
  // et le solde n'a pas bougé
  assertEquals(eur("50.00"), compte.getSolde());
}
\end{lstlisting}
    \end{column}
    \begin{column}{0.37\textwidth}
      \begin{itemize}
        \item \texttt{assertThrows} échoue si \textbf{aucune} exception (ou une exception d'un autre type) n'est levée
        \item Il renvoie l'exception : on peut l'inspecter (message, données)
        \item Vérifier aussi que l'état n'a \textbf{pas changé} après l'erreur
        \item \alert{À éviter} : \texttt{try \{ ... fail(); \} catch (...) \{\}}
      \end{itemize}
    \end{column}
  \end{columns}
\end{frame}

\begin{frame}[fragile]{Grouper les assertions : \texttt{assertAll}}
  \begin{columns}[T]
    \begin{column}{0.6\textwidth}
\begin{lstlisting}[basicstyle=\lstcode\scriptsize]
@Test
void nouveauCompte_estInitialise() {
  Compte compte = new Compte("FR76-001", eur("100.00"));

  assertAll("compte créé",
    () -> assertEquals("FR76-001", compte.getNumero()),
    () -> assertEquals(eur("100.00"), compte.getSolde()),
    () -> assertTrue(compte.getOperations().isEmpty()));
}
\end{lstlisting}
\begin{lstlisting}[style=sortie]
compte créé (2 failures)
  expected: <100.00> but was: <0.00>
  expected: <true> but was: <false>
\end{lstlisting}
    \end{column}
    \begin{column}{0.37\textwidth}
      \begin{itemize}
        \item Sans \texttt{assertAll}, le test s'arrête au \textbf{premier} échec : on corrige, on relance, on découvre le suivant\ldots
        \item Avec \texttt{assertAll}, toutes les assertions sont exécutées et \textbf{tous} les échecs sont signalés
        \item Utile pour vérifier plusieurs propriétés d'un même objet
      \end{itemize}
    \end{column}
  \end{columns}
\end{frame}

\begin{frame}[fragile]{Organiser ses tests : \texttt{@Nested} et \texttt{@DisplayName}}
  \begin{columns}[T]
    \begin{column}{0.56\textwidth}
\begin{lstlisting}[basicstyle=\lstcode\scriptsize]
@DisplayName("Un compte")
class CompteTest {
  @Nested @DisplayName("avec un solde de 100 €")
  class SoldeDe100 {
    Compte compte;

    @BeforeEach void init() {
      compte = new Compte("FR76-001", eur("100.00"));
    }
    @Test @DisplayName("accepte un débit de 30 €")
    void debit30() { ... }
    @Test @DisplayName("refuse un débit de 150 €")
    void debit150() { ... }
  }
  @Nested @DisplayName("clôturé")
  class Cloture { ... }
}
\end{lstlisting}
    \end{column}
    \begin{column}{0.42\textwidth}
      Rapport affiché par l'IDE :
\begin{lstlisting}[style=sortie]
Un compte
  avec un solde de 100 €
    ✔ accepte un débit de 30 €
    ✔ refuse un débit de 150 €
  clôturé
    ✔ refuse tout débit
    ✘ refuse tout crédit
\end{lstlisting}
      \begin{itemize}
        \item Une classe imbriquée par \textbf{contexte}
        \item Le \texttt{@BeforeEach} de chaque contexte prépare l'état
        \item Les tests se lisent comme une \textbf{spécification}
      \end{itemize}
    \end{column}
  \end{columns}
\end{frame}

\begin{frame}[fragile]{AssertJ : des assertions plus lisibles}
  \begin{columns}[T]
    \begin{column}{0.56\textwidth}
\begin{lstlisting}[basicstyle=\lstcode\scriptsize]
import static org.assertj.core.api.Assertions.*;

assertThat(compte.getSolde()).isEqualByComparingTo("70");

assertThat(compte.getOperations())
    .hasSize(2)
    .extracting(Operation::getType)
    .containsExactly(TypeOperation.CREDIT,
                     TypeOperation.DEBIT);

assertThatThrownBy(() -> compte.debiter(new BigDecimal("500")))
    .isInstanceOf(SoldeInsuffisantException.class)
    .hasMessageContaining("solde insuffisant");
\end{lstlisting}
    \end{column}
    \begin{column}{0.42\textwidth}
      \begin{itemize}
        \item Se lit de gauche à droite : « \textit{vérifie que} X \textit{est}\ldots{} »
        \item L'IDE propose les assertions adaptées au type (\texttt{.} puis complétion)
        \item Des messages d'échec bien plus clairs :
      \end{itemize}
\begin{lstlisting}[style=sortie]
assertTrue(types.contains(CREDIT));
 -> expected: <true> but was: <false>

assertThat(types).contains(CREDIT);
 -> Expecting ArrayList:
      [DEBIT]
    to contain:
      [CREDIT]
\end{lstlisting}
    \end{column}
  \end{columns}
\end{frame}

\begin{frame}[fragile]{Un piège classique : \texttt{BigDecimal} et \texttt{equals}}
  \begin{columns}[T]
    \begin{column}{0.52\textwidth}
\begin{lstlisting}[basicstyle=\lstcode\scriptsize]
Compte compte = new Compte("FR76-001",
                           new BigDecimal("100.00"));
compte.debiter(new BigDecimal("30"));

// ÉCHOUE : expected: <70> but was: <70.00>
assertEquals(new BigDecimal("70"), compte.getSolde());

// passe : compare la valeur, pas l'échelle
assertEquals(0, new BigDecimal("70")
                 .compareTo(compte.getSolde()));

// passe, et plus lisible
assertThat(compte.getSolde()).isEqualByComparingTo("70");
\end{lstlisting}
    \end{column}
    \begin{column}{0.45\textwidth}
      \begin{itemize}
        \item \texttt{BigDecimal.equals} compare la valeur \textbf{et} le nombre de décimales : \texttt{70} $\neq$ \texttt{70.00}
        \item Un test qui échoue n'indique pas toujours un bug du code : le \textbf{test} peut être faux
        \item Toujours \textbf{voir un test échouer} au moins une fois, pour être sûr qu'il teste quelque chose
      \end{itemize}
    \end{column}
  \end{columns}
\end{frame}

\rdemo{2}{Lire un échec}{%
  \item Cinq tests rouges exprès : comparer les messages d'échec
  \item \texttt{assertTrue} : « expected true but was false »\ldots{} et alors ?
  \item AssertJ : ce qu'on attendait \textbf{et} ce qu'on a obtenu
  \item \texttt{70} contre \texttt{70.00} : c'est le test qui est faux
  \item Avec et sans \texttt{assertAll} : un échec ou tous les échecs
}{\texttt{MessagesEchecTest}}

% ================================================================
\section{Organisation du module et projet fil rouge}

\begin{frame}{Six semaines : un cours, puis un TP}
  \small
  \begin{datacard}
    \begin{tabular}{@{}clll@{}}
      \toprule
      \textbf{Sem.} & \textbf{Cours (45 min)} & \textbf{TP (1 h 30)} & \textbf{Fin de semaine} \\
      \midrule
      1 & Qualité et panorama des tests & JUnit avancé sur le projet & \\
      2 & User stories, acceptation, faisabilité & Stories, grille de faisabilité, Cucumber & \\
      3 & Testabilité et doublures de test & Mockito & \alert{Quiz 1} \\
      4 & Tests d'intégration avec une base & Testcontainers, Spring, MockMvc & \\
      5 & Tests de bout en bout & Playwright, Page Object & \\
      6 & Environnements, CI, couverture & Pipeline GitLab CI, JaCoCo, PIT & \alert{Quiz 2} \\
      \bottomrule
    \end{tabular}
  \end{datacard}
  \vspace{2mm}
  \begin{center}
    Les deux quiz ont lieu dans des \textbf{séances à part} : ils ne prennent pas de temps sur les TP.
  \end{center}
\end{frame}

\begin{frame}{Le projet fil rouge : l'application bancaire}
  \begin{columns}[c]
    \begin{column}{0.6\textwidth}
      \begin{center}
      \begin{tikzpicture}[node distance=0.32cm]
        \node[layer] (front) {Front Vue 3 (navigateur)};
        \node[layer, below=of front] (ctrl) {Contrôleurs REST};
        \node[layer, below=of ctrl] (svc) {Services (\texttt{VirementService}\ldots)};
        \node[layer, below=of svc] (dom) {Domaine (\texttt{Compte}, \texttt{Client}\ldots)};
        \node[layer, below=of dom] (repo) {Repositories Spring Data};
        \node[layer, below=of repo, draw=qamuted] (db) {PostgreSQL};
        \draw[arr] (front) -- node[right, font=\tiny, text=qamuted] {HTTP / JSON} (ctrl);
        \draw[arr] (ctrl) -- (svc);
        \draw[arr] (svc) -- (dom);
        \draw[arr] (dom) -- (repo);
        \draw[arr] (repo) -- (db);
        \node[ext, right=0.9cm of svc, yshift=0.55cm] (taux) {API taux de change};
        \node[ext, right=0.9cm of svc, yshift=-0.2cm] (notif) {Notification e-mail/SMS};
        \node[ext, right=0.9cm of svc, yshift=-0.95cm] (clock) {Horloge (\texttt{Clock})};
        \draw[arr, qaaccent2] (svc.east) -- (taux.west);
        \draw[arr, qaaccent2] (svc.east) -- (notif.west);
        \draw[arr, qaaccent2] (svc.east) -- (clock.west);
        \node[tag, left=0.15cm of front] {TP5};
        \node[tag, left=0.15cm of ctrl] {TP4};
        \node[tag, left=0.15cm of svc] {TP3};
        \node[tag, left=0.15cm of dom] {TP1};
        \node[tag, left=0.15cm of repo] {TP4};
      \end{tikzpicture}
      \end{center}
    \end{column}
    \begin{column}{0.38\textwidth}
      \begin{itemize}
        \item Le domaine de R3.07 : clients, comptes, opérations, virements
        \item Projet de départ fourni : Spring Boot, Maven, docker-compose, front Vue 3 déjà écrit
        \item Collaborateurs externes (en pointillé) : à remplacer par des \textbf{doublures} dans les tests
        \item Les scénarios Cucumber (TP2) traversent toutes les couches ; la pipeline (TP6) lance tout
      \end{itemize}
    \end{column}
  \end{columns}
\end{frame}

\begin{frame}{Rendus et évaluation}
  \begin{columns}[T]
    \begin{column}{0.5\textwidth}
      \begin{block}{Rendus de TP}
        \begin{itemize}
          \item En \textbf{binôme}, un dépôt Git par binôme
          \item \textbf{Chaque} TP est rendu : tag Git \texttt{tp1}, \ldots, \texttt{tp6}, avant la date limite indiquée dans le sujet
          \item Un seul TP sera corrigé, choisi à la fin du module : \alert{tous comptent}
        \end{itemize}
      \end{block}
    \end{column}
    \begin{column}{0.46\textwidth}
      \begin{block}{Coefficient d'assiduité}
        \small
        Un rendu manquant ou en retard $=$ 1 défaut ;\\ deux rendus incomplets $=$ 1 défaut
        \vspace{1mm}
        \begin{center}
        \begin{tabular}{@{}lccc@{}}
          \toprule
          Défauts & 0 & 1 & 2 et plus \\
          Coefficient & 1 & 0,9 & 0,75 \\
          \bottomrule
        \end{tabular}
        \end{center}
      \end{block}
    \end{column}
  \end{columns}
  \vspace{3mm}
  \begin{datacard}
    \textbf{Note} $=$ moyenne\big( moyenne(quiz 1, quiz 2) , note du TP corrigé $\times$ coefficient d'assiduité \big)
  \end{datacard}
\end{frame}

\begin{frame}{À retenir}
  \small
  \begin{itemize}
    \item Un bug coûte d'autant plus cher qu'il est découvert tard : tester \textbf{tôt} et \textbf{souvent}
    \item Vérification (le produit est-il correct ?) $\neq$ validation (est-ce le bon produit ?)
    \item Quatre niveaux : unitaire, intégration, bout en bout, acceptation ; une \textbf{pyramide} : beaucoup de tests rapides, peu de tests lents
    \item Non-régression : rejouer automatiquement tous les tests à chaque modification ; un bug corrigé $=$ un test ajouté
    \item Un bon test : Arrange / Act / Assert, un nom explicite, FIRST
    \item Choisir ses cas : partitions et valeurs limites, puis tests paramétrés
  \end{itemize}
  \vspace{3mm}
  \begin{exampleblock}{TP1}
    Prise en main du projet fil rouge et tests JUnit avancés de la couche domaine.
  \end{exampleblock}
\end{frame}

\rquestionslide

\end{document}
